% !TeX root = ../../Book2.tex
% 纲要标题：### 3.1 自由模
% 纲要位置：计划纲要.md，第 887 行。

\section{自由模}
\label{b2:sec:0301}

向量空间有基，因而可以通过任意指定基向量的像来构造线性映射。对环 \(R\) 上的模也可以提出同样的问题：给定一族元素，何时能够不受约束地指定它们的像？这样的元素必须生成整个模，又不能满足任何系数不全为零的有限 \(R\)-线性关系。这正是自由模的出发点。本章沿用幺性左模的约定；讨论秩时会另外说明对环的要求。

% 纲要要点（供目录核对）：
%
% * 基与自由模的泛性质
% * 自由模的秩及其适用假设
% * 自由模的子模不必自由
%

\subsection{基与自由模的泛性质}
\label{b2:subsec:0301:01}

\begin{definition}\label{b2:def:free-module}
设 \(R\) 是环，\(M\) 是左 \(R\) 模，\((e_i)_{i\in I}\) 是其中一族元素。若每个 \(m\in M\) 都能唯一写成
\[
m=\sum_{i\in I}r_i e_i,\qquad r_i\in R,
\]
其中仅有限多个 \(r_i\) 非零，则称这族元素是 \(M\) 的一组\term[moji]{基}[basis]，称 \(M\) 为\term[ziyoumo]{自由模}[free module]。等价地，这族元素生成 \(M\)，且每个有限关系 \(\sum_i r_i e_i=0\) 都迫使所有系数为零。
\end{definition}

坐标模 \(R^{(I)}=\bigoplus_{i\in I}R\) 的标准向量 \(e_i\) 构成基。反过来，任何带有上述基的自由模，都由 \((r_i)\mapsto\sum_i r_i e_i\) 与 \(R^{(I)}\) 同构。这里使用直和，因为一般模中没有无限求和的运算；直积 \(R^I\) 的每个元素则可以有无限多个非零坐标。有限集合 \(I\) 时写作 \(R^n\)，空集合给出零模。

\begin{theorem}[自由模的泛性质]\label{b2:thm:free-module-universal}
设 \(R\) 是环，\(F\) 是有基 \((e_i)_{i\in I}\) 的左 \(R\) 模。对任意左 \(R\) 模 \(N\) 及任意一族 \((n_i)_{i\in I}\subseteq N\)，存在唯一模同态 \(f:F\to N\) 满足 \(f(e_i)=n_i\)。其公式为
\[
f\left(\sum_i r_i e_i\right)=\sum_i r_i n_i.
\]
反过来，若 \(F\) 中的一族元素对所有 \(N\) 及所有指定像都满足这一存在唯一性，则这族元素是 \(F\) 的基。
\end{theorem}
\begin{proof}
有基时，每个元素的有限坐标唯一，所以公式不依赖表达方式。相加或在左侧乘 \(r\in R\)，坐标分别相加或变为 \(rr_i\)，从公式可直接得到 \(R\) 线性。任何同态都必须把有限线性组合送到像的同样组合，因而唯一。

反过来，要证明给定的 \((e_i)\) 是基，只须把 \(F\) 与具有形式基 \((u_i)\) 的坐标模 \(E=R^{(I)}\) 同构起来，并让 \(u_i\) 对应到 \(e_i\)。两者都允许任意指定这族元素的像，因而可以分别构造两个方向的映射。由已经证明的部分，有同态 \(g:E\to F\)，\(g(u_i)=e_i\)；由假设，有同态 \(h:F\to E\)，\(h(e_i)=u_i\)。于是 \(hg\) 和 \(\operatorname{id}_E\) 在标准基上相同，由前一部分的唯一性相等；\(gh\) 和 \(\operatorname{id}_F\) 在所有 \(e_i\) 上相同，由假设的唯一性相等。因此 \(g\) 是同构，将 \(E\) 的标准基送成 \(F\) 的基。
\end{proof}

令 \(\iota:I\to F\)、\(\iota(i)=e_i\)，并令 \(\nu:I\to N\)、\(\nu(i)=n_i\)。图按底集合理解，两条实线是集合映射，虚线是唯一的左 \(R\) 模同态。
\[
\begin{tikzcd}[column sep=3em,row sep=2.5em]
I\arrow[r,"\iota"]\arrow[dr,"\nu"']&F\arrow[d,dashed,"\exists!\,f"]\\
&N.
\end{tikzcd}
\]
例如，\(\mathbb Z^2\) 中的 \((1,0),(0,1)\) 可以任意映到一个交换群的两个元素；但 \(\mathbb Z/6\mathbb Z\) 的生成元 \(\bar1\) 只能映到满足 \(6n=0\) 的元素。限制来自生成元的关系，并非来自是否只用一个生成元。若 \(R\ne0\)，自由模中的任何基向量都没有非零标量将它零化，因此非零自由模是忠实的。

\begin{proposition}\label{b2:prop:free-lifting}
设 \(R\) 是环，\(F,M,N\) 是左 \(R\) 模，且 \(F\) 自由。若 \(q:M\to N\) 是满同态，则每个模同态 \(f:F\to N\) 都可以写成 \(qg\)，其中 \(g:F\to M\) 为模同态。
\end{proposition}
\begin{proof}
所求映射 \(g\) 把 \(f\) 沿满同态 \(q\) 提升到 \(M\)，使下面的三角形交换；其中箭头均为左 \(R\) 模同态，虚线表示要构造的映射。
\[
\begin{tikzcd}[column sep=3em,row sep=2.5em]
&M\arrow[d,two heads,"q"]\\
F\arrow[r,"f"']\arrow[ur,dashed,"g"]&N.
\end{tikzcd}
\]
取 \(F\) 的基 \((e_i)\)。由满射性，逐个选取 \(m_i\in M\) 使 \(q(m_i)=f(e_i)\)。自由模的泛性质给出唯一满足 \(g(e_i)=m_i\) 的同态 \(g\)。\(qg\) 与 \(f\) 在基上相同，所以相等。基无限时，选取这族原像使用本书采用的选择公理；有限基时只作有限次选取。
\end{proof}

这个提升通常不唯一，因为每个 \(m_i\) 都可以加上 \(\ker q\) 中的元素。唯一的是选定这些原像后得到的同态。自由模的这一性质将在分裂短正合列中直接使用。

\subsection{自由模的秩与适用条件}
\label{b2:subsec:0301:02}

自由模一旦选定基，基的指标集当然有基数；要把这个数称为模本身的秩，还须证明它不随基改变。一般环上，这一步不能照搬向量空间的消元证明，因为非零系数未必可逆。

\begin{theorem}[交换环上自由模的秩不变性]\label{b2:thm:free-rank-invariance}
设 \(R\) 是非零交换环，\(I,J\) 是集合。如果 \(R^{(I)}\cong R^{(J)}\)，则 \(|I|=|J|\)。因此，若自由 \(R\) 模 \(F\) 有一组基，就把其基数称为 \(F\) 的\term[ziyoumozhichi]{秩}[rank]，记为 \(\operatorname{rank}_R F\)。
\end{theorem}
\symidx{\(\operatorname{rank}_RF\)：自由模的秩}
\begin{proof}
我们把自由模的所有系数同时模掉一个极大理想，使系数环变成域，再用向量空间的维数不变性。由于 \(R\) 是非零交换环，\((0)\) 是真理想，第一卷第十章命题 10.2.7 保证它包含在某个极大理想 \(\mathfrak m\) 中；该章命题 10.2.6 保证商环 \(k=R/\mathfrak m\) 是域。

记 \(\mathfrak mF\) 为所有有限和 \(\sum a_\nu x_\nu\) 组成的子模，其中 \(a_\nu\in\mathfrak m\)、\(x_\nu\in F\)。商模 \(F/\mathfrak mF\) 上的 \(R\) 作用因 \(\mathfrak m\) 作用为零而下降为 \(k\) 的作用，所以它是 \(k\) 向量空间。对带基的自由模而言，取这个商就是把每个坐标系数 \(r_i\) 换成 \(r_i+\mathfrak m\)，原基向量的像成为域 \(k\) 上的基。具体地，当 \(F=R^{(I)}\) 时，逐坐标取剩余类给出映射
\[
R^{(I)}\longrightarrow k^{(I)},\qquad
(r_i)_{i\in I}\longmapsto(r_i+\mathfrak m)_{i\in I}.
\]
每个有限支集剩余类向量都有原像，而此映射的核恰为由坐标都在 \(\mathfrak m\) 中的有限支集向量组成的 \(\mathfrak mR^{(I)}\)。因此
\[
R^{(I)}/\mathfrak mR^{(I)}\cong k^{(I)}.
\]
模同构 \(f:F\to G\) 将 \(\mathfrak mF\) 映成 \(\mathfrak mG\)：正向包含来自 \(f(ax)=a\cdot f(x)\)，反向包含对 \(f^{-1}\) 使用同一公式。于是它诱导 \(k^{(I)}\cong k^{(J)}\)。由\cref{b2:thm:vector-dimension}的向量空间维数不变性，\(|I|=|J|\)。
\end{proof}

对除环也有基数不变性，有限情形可以用单侧消元证明；上面的极大理想方法则专门利用交换环的域商。零环必须排除：这时任意 \(R^{(I)}\) 都是零模，指标集的大小不能由模恢复。

\ProseParagraphBreak
非交换环甚至可能有 \(R\cong R^2\)。设 \(V\) 有可数基 \(v_0,v_1,\ldots\)，令 \(R=\operatorname{End}_k(V)\)，乘法为复合。定义
\[
s_0(v_i)=v_{2i},\qquad s_1(v_i)=v_{2i+1},
\]
并令 \(p_0\) 把偶数位置基向量 \(v_{2i}\) 送到 \(v_i\)，将奇数位置送到零；\(p_1\) 对奇数位置作相应操作。它们满足
\[
p_i s_j=\delta_{ij}\operatorname{id}_V,\qquad
s_0p_0+s_1p_1=\operatorname{id}_V.
\]
把基分成偶数位置与奇数位置，就是把输入空间写成
\[
V=\operatorname{span}_k\{v_{2i}:i\geq0\}
\oplus\operatorname{span}_k\{v_{2i+1}:i\geq0\}.
\]
一个算子 \(c:V\to V\) 由它在这两部分输入上的作用共同决定。\(s_0,s_1\) 把整个 \(V\) 分别识别为这两部分，所以两份限制可重新写成 \(V\to V\) 的算子 \(cs_0,cs_1\)：对 \(v,w\in V\)，
\[
c(s_0v+s_1w)=(cs_0)(v)+(cs_1)(w).
\]
反过来，给定两个算子 \(a,b\)，先用 \(p_0,p_1\) 拆出输入的两部分，再分别施加 \(a,b\) 并相加，就得到 \(ap_0+bp_1\)。因此定义
\[
\Phi:R^2\longrightarrow R,\quad(a,b)\longmapsto ap_0+bp_1,
\]
\[
\Psi:R\longrightarrow R^2,\quad c\longmapsto(cs_0,cs_1).
\]
它们确为左 \(R\) 模同态：对 \(r\in R\)，有 \(\Phi(ra,rb)=r\Phi(a,b)\)、\(\Psi(rc)=r\Psi(c)\)。固定算子乘在右侧，是因为复合时先拆分输入，而左侧的 \(r\) 仍统一作用于输出。上面的算子恒等式又给出
\[
\begin{aligned}
\Psi\Phi(a,b)&=\bigl((ap_0+bp_1)s_0,(ap_0+bp_1)s_1\bigr)=(a,b),\\
\Phi\Psi(c)&=c(s_0p_0+s_1p_1)=c.
\end{aligned}
\]
所以这两映射互逆。这个环的自由模因此可以有基数不同的基；在一般环上没有额外假设时，应说“具有 \(n\) 元基”，不能把 \(n\) 默认为模的一个不变量。但这并不表明所有非交换环都失去秩不变性，前面提到的除环就仍有这一性质。交换环的结论与此例可对照 Roman\cite[Chapter 5, Theorem 5.1; Example 5.1, pp. 129--130]{Roman2008AdvancedLinearAlgebra}。

\subsection{自由模的非自由子模}
\label{b2:subsec:0301:03}

向量空间的子空间仍有基，一般环上却不成立。先看一个由零因子造成的例子。取 \(R=k[\varepsilon]/(\varepsilon^2)\)，其中 \(\varepsilon\ne0\)，左正则模 \(R\) 自由，子模 \(I=(\varepsilon)\) 非零且满足 \(\varepsilon I=0\)。如果 \(I\) 自由，它至少有一个基向量 \(e\)，关系 \(\varepsilon e=0\) 就迫使 \(\varepsilon=0\)，矛盾。

\ProseParagraphBreak
去掉零因子仍不足够。令 \(R=k[x,y]\)，\(I=(x,y)\)。若 \(I\) 有两个不同的基向量 \(e_1,e_2\)，它们也是 \(R\) 中的非零多项式。环的交换性给出关系
\[
e_2\cdot e_1+(-e_1)\cdot e_2=0.
\]
在这条模关系中，作为基向量的是 \(e_1,e_2\)，作为环中系数的是 \(e_2,-e_1\)；两个系数均非零，与基的线性无关性矛盾。因此 \(I\) 若自由，基至多有一个元素；由于 \(I\ne0\)，它必须是主理想 \((h)\)。但 \(h\) 同时整除 \(x\) 和 \(y\)：若 \(h\) 非常数，由全次数的乘法公式及 \(x=ha\)，\(a\) 是非零常数，\(h\) 是 \(x\) 的常数倍，于是不能整除 \(y\)。所以 \(h\) 必须可逆，这又意味着 \(I=R\)，与 \(I\) 中所有多项式的常数项都为零矛盾。因此 \(I\) 是自由模 \(R\) 的非自由子模。

\ProseParagraphBreak
还要区分“子模自由”与“子模有补模”。\(2\mathbb Z\) 以 \(2\) 为基，是自由 \(\mathbb Z\) 模，却不是 \(\mathbb Z\) 的直和因子。若 \(\mathbb Z=2\mathbb Z\oplus N\)，则对 \(n\in N\)，有 \(2n\in2\mathbb Z\cap N=0\)，由整数的无挠性得 \(n=0\)，这又不能生成 \(\mathbb Z\)。第四章将证明主理想整环上有限秩自由模的子模仍自由；它保证关系可以有限地整理，却并不保证每个子模都有补模。
